%==========================================================================
\section{Static semantics}
\subsection*{New Type}
ReflectInfo is a synthetic class, like Input and Output. It has these
built-in methods, called with ordinary method-call syntax:

\begin{center}\small
\begin{tabularx}{\linewidth}{|*{5}{>{\raggedright\arraybackslash\ttfamily\footnotesize}X|}}
\hline
String class\_name() & String parent\_name() & int field\_count() &
  String field\_name(int i) & String field\_type(int i) \\
\hline
int method\_count() & String method\_name(int i) & String method\_return(int i) &
  int method\_arity(int i) & String method\_param\_type(int i, int j) \\
\hline
\end{tabularx}
\end{center}

ReflectInfo values can be stored, passed and returned like any class
reference, including null. A program cannot declare a class named ReflectInfo
(\code{E\_RESERVED\_CLASS\_NAME}), extend it, or construct it with new
(\code{E\_REFLECT\_RESERVED\_CLASS}).

\subsection*{Subtyping}
No new subtype relations. ReflectInfo is unrelated to every user class, and
null stays compatible with every class type, ReflectInfo included.

\subsection*{Scoping and name resolution}
No scoping changes. The new keywords are reserved. The member-name operand is
an ordinary String expression. The type checker does not resolve it, so it
never reports \code{E\_UNKNOWN\_METHOD} or \code{E\_UNKNOWN\_VARIABLE} for it,
and lookup happens at runtime. For the same reason, invoke's and
new\_instance's argument counts and types, and whether the requested type
matches the member's declared type, are checked at runtime, not statically.

\subsection*{Encapsulation}
In LO-4, a field is reachable only by its bare name
inside its own class's methods and constructors, i.e.\ only on
\code{this}, whereas getfield and setfield reach any field of any object the caller
holds a reference to, so reflection grants cross-object field access that
LO-4 lacks. This weakens LO-4's encapsulation guarantee, but keeps the
implementation simple, requiring no per-class reflectability flag and no extra
runtime check. However, every such access is syntactically
visible through the getfield or setfield keyword and remains type-checked.

\subsection*{Type Rules}
\begin{center}
\renewcommand{\arraystretch}{2.6}
\begin{tabular}{@{}l c@{}}
\toprule
(T-Reflect) &
  $\infrule{\Gamma \vdash e : C \quad C \text{ is a user class}}
           {\Gamma \vdash \tm{reflect}(e) : \tm{ReflectInfo}}$ \\
(T-GetField) &
  $\infrule{\Gamma \vdash e : C \quad C \text{ is a user class} \quad \Gamma \vdash n : \tm{String} \quad T \neq \tm{void}}
           {\Gamma \vdash \tm{getfield}(T, e, n) : T}$ \\
(T-SetField) &
  $\infrule{\Gamma \vdash \tm{getfield}(T, e, n) : T \quad \Gamma \vdash v : S \quad S \le T}
           {\Gamma \vdash \tm{setfield}(T, e, n, v)\ \text{ok}}$ \\
(T-Invoke) &
  $\infrule{\Gamma \vdash e : C \quad C \text{ is a user class} \quad \Gamma \vdash n : \tm{String} \quad \Gamma \vdash a_i : A_i \quad T \neq \tm{void}}
           {\Gamma \vdash \tm{invoke}(T, e, n, a_1, \ldots, a_k) : T}$ \\
(T-InvokeStmt) &
  $\infrule{\Gamma \vdash e : C \quad C \text{ is a user class} \quad \Gamma \vdash n : \tm{String} \quad \Gamma \vdash a_i : A_i \quad T = \tm{void}}
           {\Gamma \vdash \tm{invoke}(T, e, n, a_1, \ldots, a_k);\ \text{ok}}$ \\
(T-NewInstance) &
  $\infrule{\Gamma \vdash n : \tm{String} \quad \Gamma \vdash a_i : A_i \quad T \text{ is a user class}}
           {\Gamma \vdash \tm{new\_instance}(T, n, a_1, \ldots, a_k) : T}$ \\
(T-Allocate) &
  $\infrule{\Gamma \vdash n : \tm{String} \quad T \text{ is a user class}}
           {\Gamma \vdash \tm{allocate}(T, n) : T}$ \\[1ex]
\bottomrule
\end{tabular}
\end{center}

\subsection*{Diagnostics}
\begin{center}\small
\begin{tabularx}{\linewidth}{@{}l X@{}}
\toprule
\textbf{Error Code} & \textbf{Trigger Condition} \\
\midrule
\code{E\_REFLECT\_RECEIVER\_TYPE} & Receiver operand does not evaluate to a user class reference type. \\
\code{E\_REFLECT\_TARGET\_TYPE} & getfield or setfield requests type void, or new\_instance's or allocate's T is not a user class. \\
\code{E\_REFLECT\_VALUE\_TYPE} & The setfield value is not assignment-compatible with T. \\
\code{E\_REFLECT\_RESERVED\_CLASS} & \code{new ReflectInfo(...)} or \code{extends ReflectInfo}. \\
Existing codes & A non-String name (\code{E\_TYPE\_MISMATCH}), an unknown class in T
  (\code{E\_UNKNOWN\_CLASS}), the wrong void/non-void form of invoke
  (\code{E\_VOID\_CALL\_IN\_EXPRESSION} / \code{E\_NONVOID\_CALL\_AS\_STATEMENT}) \\
\bottomrule
\end{tabularx}
\end{center}

\noindent
\begin{minipage}[t]{0.49\linewidth}
\textbf{Well-Formed}
\begin{lstlisting}
class Box (int count;) { }
class Main () {
  int main() {
    Box b; b = new Box(2);
    setfield(int, b, "count", 4);
    return getfield(int, b, "count");
  }
}
\end{lstlisting}
\end{minipage}\hfill
\begin{minipage}[t]{0.49\linewidth}
\textbf{Ill-Formed}
\begin{lstlisting}
class Box (int count;) { }
class Main () {
  int main() {
    Box b; b = new Box(2);
    setfield(int, b, "count", true);
    return 0;
  }
}
\end{lstlisting}
\end{minipage}

This program violates the setfield rule ($S \le T$). Diagnostic:
\code{E\_REFLECT\_VALUE\_TYPE: setfield value has type bool, expected int}

